\section{Summary and Future Directions}
\label{sec:rl_summary}

This chapter presented two contributions for enabling autonomous agents to adapt to novelty in open-world environments. The \RAPidLearn{} framework formalizes the integration of symbolic planning and reinforcement learning through the concept of integrated planning tasks, defines the executor discovery problem as the core challenge of novelty recovery, and introduces knowledge-guided exploration strategies that exploit domain knowledge to achieve sample-efficient adaptation. The neurosymbolic cognitive architecture embeds \RAPidLearn{} within a broader system that connects detection, characterization, and accommodation within the evaluated symbolic task interface by combining symbolic inference, neural inference, and reinforcement learning in a cascading recovery pipeline.

The experimental evaluations demonstrate three key findings. First, \RAPidLearn{} achieves one to two orders of magnitude faster adaptation than transfer learning baselines by decomposing the global recovery problem into focused subproblems defined by failed operators. Second, knowledge-guided exploration provides measurable benefits for complex novelties that require the agent to interact with novel entities or discover spatial relations not present in its symbolic model. Third, the integrated architecture achieves robust performance across 288 novelties spanning 10 categories, including concealed novelties not anticipated during system development, demonstrating the value of combining symbolic reasoning with learning-based recovery across diverse changes in the evaluated Polycraft setting.

These results instantiate the accommodation side of the dissertation's central claim. Action abstractions expose the operator or executor whose expected behavior no longer matches the environment, allowing the architecture to revise that structure or learn a replacement executor without discarding unaffected competence. The agent thereby expands its behavioral repertoire at the point of failure rather than relearning the entire task. The accommodation demonstrated here is localized and behaviorally grounded, but scoped by the symbolic state interface, operator vocabulary, and exploration mechanisms available to the architecture.

Several limitations of the current approach point to important directions for future work.

\paragraph{Failure-driven adaptation.} As observed in \cref{ch:benchmarks}, the failure-driven nature of \RAPidLearn{} means that beneficial novelties that do not cause execution failures go undetected and unexploited. The chest novelty in the \NovelGym{} evaluation (\cref{sec:ng_results}) illustrates this clearly: hybrid agents never discovered the shortcut because their plans continued to succeed. Extending the framework to proactively detect and exploit beneficial opportunities---for example, through periodic re-planning with updated models---would address this blind spot.

\paragraph{Fixed goal assumption.} The current framework assumes a fixed goal state throughout the adaptation process. In realistic open-world settings, goals themselves may need to be revised when novelty renders them unachievable or when new opportunities arise. Integrating goal reasoning~\citep{aha2018goalreasoning} with the executor discovery mechanism would enable more flexible responses to environmental changes.

\paragraph{Exploration cost in physical environments.} The reinforcement learning component of \RAPidLearn{} requires trial-and-error interaction with the environment, which is acceptable in simulation but potentially costly or dangerous in physical robotic settings. Reducing the number of exploratory interactions---through model-based reinforcement learning, sim-to-real transfer, or more informed exploration strategies---is essential for deploying these methods on real robots. A complementary strategy is to design representations that allow previously learned behavior to remain applicable across relevant environmental and embodiment variation, thereby reducing how often online recovery learning is needed.

\paragraph{Symbolic-to-subsymbolic grounding.} The current architecture relies on a domain interface that provides pre-processed symbolic state descriptions. In real-world settings, the agent would need to ground these descriptions in raw sensory data---a challenging perception problem that interacts with novelty detection (e.g., visually recognizing a novel object and inferring its symbolic type). Tighter integration of the visual novelty detection capabilities with the symbolic reasoning components would enable end-to-end novelty handling from raw observations.

\rev{\paragraph{Restricted symbolic search spaces.} To keep exploration tractable within the task time-out, the symbolic operator-discovery strategies are deliberately restricted---for example, operator variation is applied only to single-parameter operators, and precondition discovery focuses on a small set of prominent predicates. These heuristics create blind spots: the ``must break tree tap to get rubber'' and ``cannot collect rubber from birch trees'' novelties (\cref{sec:rl_eval_architecture}) are both unsolvable by the Base agent for precisely this reason, and are recovered only when the executor learner is available. Learning to \emph{allocate} exploration effort adaptively, rather than relying on fixed domain-specific heuristics, would reduce these blind spots without incurring the full cost of unrestricted symbolic search.}

\rev{\paragraph{Purely visual novelty detection.} The integrated architecture includes a vision model for completeness, but purely visual detection is not a focus of this dissertation. Visual detections cannot drive accommodation without symbolic descriptions of what changed, and the false-positive rate remains too high for reliable task-oriented use in the Polycraft evaluation.}

\paragraph{Coordination between adaptive components.} The systemic evaluation revealed that the interaction between the agent model and the executor learner can be counterproductive: the learner's exploration triggers false novelty detections from the agent model, leading to unnecessary exploration. Developing mechanisms for coordinating multiple adaptive components---for example, by suppressing novelty detection during active learning phases, or by sharing information about the source of environmental changes between components---is an important engineering and research challenge for complex novelty-handling architectures.

Taken together, the results in this chapter show how action abstractions can turn execution failure into localized accommodation: when existing behavior is insufficient, the agent expands its repertoire at the failed operator--executor interface while preserving competence that remains valid. The next two chapters examine the complementary role of interaction abstractions. They ask when object, physical, and embodiment variation can instead be assimilated by an existing policy, avoiding recovery learning for the evaluated classes of variation.
